\section{Threats to Validity}
\label{sec:threats}

\textbf{Conclusion validity.} The 2{,}700-row evaluation provides
narrow but adequate statistical power for the cell-mean comparisons.
H3's Spearman is computed on three points and is exact ($\rho = 1.0$
because the three fire fractions are strictly ordered). H4's exact
$0.000$ delta between cap\_aware and gated at K=200 is a structural
identity rather than a statistical comparison: the two strategies
make the same decision at every window, so their cell means coincide
to numerical precision.

\textbf{Internal validity.}

\textit{(i) The $\tau_K$ vector is pre-registered, not tuned.} The
choices $\{0.20, 0.05, 0.02\}$ were locked before any
capacity-aware-detector P@50 was inspected. They were chosen by
inspection of Paper~11's calibration-seed change-point decision log
(an artifact released with Paper~11) and operational intuition.
Other $\tau_K$ vectors could shift the K=100 intervention rate but
would not change the K=50 and K=200 collapses.

\textit{(ii) Selection-policy coupling.} Inherited from Papers~7--11.
Fleet evolution is driven by HygienePrio-full under fixed weights;
all six strategies share that trajectory. Paper~9~\cite{paper9selftraj}
addressed the cross-method trajectory question; the present paper
inherits the fixed-trajectory convention.

\textit{(iii) Single $\lambda$ value.} The sweep holds $\lambda = 3$.
Different arrival rates would shift the per-window $|\Delta_w|$
distribution and re-calibrate the $\tau_K$ vector accordingly. The
qualitative finding --- a per-K threshold reaches feasibility by
approximating the static gate --- is more robust than the specific
$\tau_K$ choice.

\textbf{Construct validity.} The H1 K=50 result holds with equality
because cap\_aware reproduces lag-1 exactly when the detector never
fires. This is the substantive scientific claim, not an artifact: the
$\tau_{50} = 0.20$ choice was made to suppress fires, and the
suppression was achieved. The K=200 H4 equality is similarly substantive.

\textbf{External validity.} Synthetic EEHDA structural priors are
inherited from Papers~4--11. Real fleets with substantially different
per-window calibration-target distributions would re-locate the
feasibility region along the $\tau_K$ axis. The qualitative result
(cap\_aware $\to$ gated at the two capacity extremes) is robust to
distributional shifts within the studied calibration-grid family.
